\documentclass[11pt,fleqn]{report}
\usepackage{graphicx}
\usepackage{amssymb}
\usepackage{longtable}
\usepackage[T1]{fontenc}

\usepackage{cml}
\lstnewenvironment{codeXML}
    {\lstset{basicstyle=\small\ttfamily,
             language=XML,
             xleftmargin=0.0in,
             xrightmargin=0.0in,
             backgroundcolor=\color{Orchid!60}, % 60% orchid, 40% white
             escapeinside={\%*}{*)}, %  .tex file characters inside %*...*) are
                                     %  treated as Tex commands.
            }
    }
    { }

\newcommand{\ModelName}{Enhanced Logging}
\newcommand{\ModelAuthor}{Gary Turner}
\newcommand{\ModelAffiliation}{Odyssey Space Research}
\newcommand{\ModelDate}{February 2024}
\pagestyle{plain}
\usepackage[parfill]{parskip}
\usepackage[margin=1.5cm]{caption}
\usepackage{subfig}

\begin{document}
\titlePage
\clearpage % Reset page number
\pagenumbering{roman}
\tableofcontents % Print table of contents
\vfill
{\Large \textbf {Revision History}}
{\large
\begin{center}
\begin{tabular}{|l||l|l|l|}
\hline
 \textbf{Version} & \textbf{Date} & \textbf{Author} & \textbf{Purpose} \\ \hline
 \hline
1 & February 2024 & Gary Turner & Initial Version \\ \hline
\end{tabular}
\end{center}
}

\chapter{Introduction} % Make a "chapter" heading
\pagenumbering{arabic} % Start text with arabic 1
The CML Enhanced Logging model is derived from the logging capabilities
provided by the Data-Collection component of the GNC-PAR package, developed for
the Artemis program -- specifically for the ANTARES and OrionSim simulations.
The GNC-PAR Data-Collection package was developed tp provide additional
logging capabilities -- beyond those provided by the Trick simulation engine
-- that would be desirable to support the process of verifying the Orion
vehicle performance.

During the verification process of this GNC-PAR package, these logging
capabilities that had built up over many years were consolidated into this
new package; this effort included:

\begin{itemize}
  \item thorough review of the existing code, updating as needed and
        writing within a consistent framework;
  \item identification of the implicit requirements that had led to the
        development of the package capabilities;
  \item formalization and documentation of those implicit requirements to
        make them explicit;
  \item verification of the internal data processing capabilities;
  \item overall model documentation;
  \item relocation of the updated model from an Orion-specific location
        (GNC-PAR) to the EG-wide CML library.
\end{itemize}
Two closely-related but independent packages were created in this process:

\begin{itemize}
  \item \textbf{Enhanced Logging} provides the interface to use the enhanced
        logging capabilities while generating continuous data files.
  \item \textbf{Summary Logging} provides the interface to generate only
        summary logs of data created using the enhanced logging
        capabilities.
        See \textit{CML: models/utilities/summary\_logging} for details.
\end{itemize}

This document describes the former of these packages.

\chapter {Requirements}
\section{Nomenclature and Concepts}\label{sec:2:concepts}

The following concepts are drawn from the CML \textit{Compound-Event} model
upon which this model is built:
(see \textit{CML: vehicle\_management/compound\_events\_manager})
\begin{itemize}
  \item \textbf{Action Triggers:}
        An \textit{event} may include a set of one or more
        \textit{action-triggers}; when these triggers
        are collectively satisfied, the event is considered \textit{triggered}.
  \item \textbf{Arm / Arming Triggers:}
        An \textit{event} may include an arming phase, which provides
        a set of one or more  \textit{arming-triggers}; when these triggers
        are collectively satisfied, the event is considered armed.
        After an event is armed, or if it does not have an arming phase, it
        will start monitoring its \textit{action-triggers} and
        \textit{disarming-triggers}.
  \item \textbf{Disarm / Disarming Triggers:}
        An \textit{event} may include a disarming option, which provides
        a set of one or more  \textit{disarming-triggers}. These triggers
        are monitored after the \textit{arming-triggers} have been satisfied
        and the event consequently in an \textit{armed} state. When the
        \textit{disarming-triggers}
        are collectively satisfied, the  \textit{armed} status of the
        associated event is removed.
        After an event is disarmed, the arming phase is re-entered, and the
        arming requirements must be satisfied again before the event starts
        monitoring its  \textit{action-triggers} again.
  \item \textbf{Event:}
        An event is an encapsulation of arming-triggers, disarming-triggers,
        and action-triggers that, collectively, lead to the prescribed
        actions in response to the conditions of those triggers being
        satisfied. Available actions that are built-in to each event
        include (see CML Events-manager \textit{CML:
        models/vehicle\_management/events\_manager} for more details):
  \begin{itemize}
    \item Setting of a boolean flag \textit{event\_triggered} to true. All
          events set this flag; it is commonly used by other models as an
          indication of the completion of some event.
    \item Posting of a message to standard-out.
    \item Assignment of values to other variables elsewhere in the
          simulation.
    \item Add or remove subscriptions to/from other models in the
          simulation, potentially affecting their \textit{active} status.
    \item Disable other models elsewhere in the simulation, forcing their
          active status to false regardless of their use by other models.
    \item Execute additional methods to apply additional algorithms. These
          additional algorithms are defined outside the structure of the
          event; the event provides only the interface to them making them
          effectively limitless in capability.
  \end{itemize}
  \item \textbf{Trigger (noun):} A trigger is a logical comparison that is
        used to define event conditions and consequently identify when event
        conditions have been satisfied.
  \item \textbf{Trigger (verb):} An object is considered \textit{triggered}
        when its logical tests have been satisfied. An event may include
        multiple triggers, compounded with AND and OR logic in multiple ways.
\end{itemize}

The following concepts are introduced for this model. These concepts are
interrelated and can only be truly understood as a collective. There isn't
really an obvious starting point so they are presented in alphabetical
order, with concepts identified within each description by italicized text.
\begin{itemize}
  \item \textbf{Condition:} Each \textit{logging-group} has a set of
        \textit{conditions}, each of which has an effective true/false status.
        If a \textit{condition} is true, the
        current values of the group's \textit{simulation-variables} will be
        \textit{processed} to determine whether to  \textit{record} any of the
        current values of the group's \textit{simulation-variables} in the
        group's \textit{logging-variables}.
  \item \textbf{Log:} Logging is the explicit writing of a set of values
        contained within each logging-variable to an output files.
        When logging, the model may be configured to log the contents of each
        \textit{logging-group} into its own output file, or to consolidate
        all contents into one file.
        There is no finer division for logging than the \textit{logging-group}.
  \item \textbf{Logging-group:} A logging group encapsulates a set of
        \textit{logging-variables} and a linked set of \textit{conditions}. Each
        \textit{logging-variable} maintains space for a \textit{recorded}
        value for each of the  \textit{conditions} within the logging-group.
  \item \textbf{Logging-variables:} Logging-variables are a component of the
        model. Each \textit{logging-variable} is associated with a single
        \textit{simulation-variable}. Each maintains a set of
        \textit{recorded} values based on the values of the associated
        \textit{simulation-variable} at the time(s) that each of the associated
        \textit{conditiosn} were satisfied.
        Each \textit{condition} in the \textit{logging-group} has a dedicated
        slot for an associated \textit{recorded} value in each of the
        \textit{logging-variable}'s set of recorded values.
  \item \textbf{Processing a variable:} When one or more \textit{conditions}
        are satsified, the values of the \textit{simulation-variables} will be
        \textit{processed}. The algorithm applied is determined by the
        \textit{record-type} of the satisfied condition(s) and includes
        consideration of the current value \textit{recorded} by the
        corresponding
        \textit{logging-variable} in the slot(s) reserved for the satisfied
        \textit{condition(s)}.
  \item \textbf{Record:} Recording is the internal storage of a value from
        a simulation-variable into one of the slots in the corresponding
        logging-variable. The associated condition that led to the
        recording determines which slot to use.
  \item \textbf{Record-type:} The record-type is a configuration setting in
        each \textit{condition} that determines how each
        \textit{simulation-variable's} value should be
        \textit{processed} when the \textit{condition} is satisfied.
  \begin{itemize}
    \item For example, one of the record-types (see \ref{sec:3:FIXME}
          for a complete list) instructs the model to \textit{record}
          the maximum
          value of a variable. When this condition is satisfied, the values
          of each simulation-variable in the same group are compared against
          the \textit{recorded} values of each logged-variable for this specific
          condition. For those variables for which the new values of the
          simulation-variables are
          greater than the corresponding recorded values, the recorded
          values are replaced. New values less than (or equal to) the
          corresponding recorded values are ignored.
  \end{itemize}
  \item \textbf{Simulation-variables:} Simulation-variables are those variables
        found in the simulation; the overall goal of this model is to
        \textit{log} a selected set of data based on values that the
        \textit{simulation-variables} have
        been \textit{recorded} as having at some point.
\end{itemize}

\section{Requirements}

TBD

\chapter {Model Specification}
See section \ref{sec:2:concepts} for fundamental concepts and nomenclature.
These will be used throughout this section without further clarification.

\section {Architectural Considerations}

\subsection {Existing and Additional Capabilities}
This model is designed to supplement the logging capabilities already
provided by the Trick simulation engine, and the CML model
\textit{Trick-logging} that provides a management-interface to control those
logging capabilities. As such, any logging operations that would be
considered routine or regularly scheduled are not duplicated by this model.

The CML model \textit{Summary-logging} provides highly capable summary data
but is limited to creating reports at simulation termination. The model
described hereinn uses couples many of thosse same capabilities with the
underlying concepts from Trick-logging to create logs of evolving data (like
Trick-logging) with the enhanced data processing capabilities found in
Summary-logging.

Within Trick-logging, there is a central “logging manager” which contains a
set of logging groups, with each group specifying a set of variables for
logging.
Each group executes at its own user-configurable frequency and produces its own
output data file, which can be configured as an ASCII (CSV) file, a binary file
using Trick’s native binary format (TRK), or binary file using the general
purpose \textit{Hierarchical Data Format 5} (HDF5) format.
Groups typically produce a complete log of all variables within
the group at the specified frequency. It is also possible to configure a group
to log only when the value of some watched variable changes.

The CML-Trick-logging extension to Trick logging (see
\textit{cml/models/utilities/trick-logging}),
already provides the interface to easily manipulate specialty logging
configurations, including:

\begin{itemize}
  \item trigger events-based logging (“log now” events) for specific groups,
        or all groups;
  \item activation and deactivation of specific groups in response to current
        simulation state or simulation events.
  \item modification of the logging frequency for specific groups, or for all
        groups, again in response to simulation state or events.
\end{itemize}

The additional capabilities provided by this model include:
\begin{itemize}
  \item Regular logging configurations of traditional variables:
  \begin{itemize}
    \item The ability is added to instantiate and configure multiple
          logging managers. One of the shortcomings of the Trick-logging
          manager is the difficulty in applying extensive reconfigurations
          at major simulation boundaries. With the ability to handle multiple
          logging managers, large reconfigurations could be implemented as
          wholesale swap-out between pre-configured managers used to control
          the logging for a specific flight phase.
    \begin{itemize}
      \item Note – it would be feasible to run multiple managers concurrently,
            and may be desirable to do so during brief periods of “hand-over”.
            However, it is not the primary intent of this capability to
            support long-term concurrent use of multiple managers, and
            configuring multiple managers to run concurrently is
            probably more challenging than utilizing other configuration
            options within a single manager.
    \end{itemize}
    \item The ability to add and remove variables from a group during
          initial configuration. Typically with Trick logging, the logging
          specifications are handled via SWIG and there is no access to a
          data-record group after it has been configured. With this model,
          data record groups can be instantiated and configured independently
          of the logging manager, and even if they are instantiated from
          within the logging manager, they can be accessed after initial
          configuration. This allows a default data recording group to be
          loaded, then modified before use.
    \begin{itemize}
      \item Note – modifications to the contents of a data record group
            can only be made up to the point at which the data record file
            has been opened.  This ensures that the contents of the data
            file are consistent throughout its population.
    \end{itemize}
  \end{itemize}
  \item Logging of STL-strings. Trick logging can have issues with
        logging STL string content. The new model is able to log STL
        string values.
  \item Application of conditional logging.
        Using the CML-Trick-logging interface, events had to be added
        to the simulation events framework to bring about changes to
        the logging configuration, or to trigger “log-now” effects.
        This model still supports that capability, but also adds a
        simple internal mechanism to support easy configuration for
        tracking conditional logging events. Because the new
        conditionals-configuration is limited to managing logging,
        it is less capable than a full events-management capability
        and so is easier to configure. Examples of the types of
        conditionals that are supported follow:
  \begin{itemize}
    \item log the contents of a group whenever some variable changes value
    \item log the contents of a group whenever some variable exceeds some
          threshold
    \item log the contents of a group whenever some variable reaches a new
          maximum value.
  \end{itemize}
        These conditionals can be combined with either OR, or AND logic,
        but the model does not support compound logic (at this time).
  \begin{itemize}
    \item A AND B AND C is supported.
    \item A OR B OR C is supported.
    \item A AND (B OR C) is not supported.
  \end{itemize}
  \item Summary logging. In many situations, it is not really necessary
        to know the current value of a variable, but instead to know the average
        value or maximum value, or similar. With Trick logging, these are
        typically acquired by either: adding a new variable and logic into
        the model itself to generate the summary value, or logging all values
        and post-processing the results to obtain the desired summary value
        from the logged raw values. This model provides logic to create a
        logging-group to record a summary including the average, or maximum,
        or minimum, etc. of each of the variables within the group.
  \begin{itemize}
    \item Note – the specification of whether to log average, maximum,
          minimum, etc. is a group-level specification, and ALL variables
          within that group will have a corresponding value logged.
          It was considered to make this a variable-level specification,
          but the potential for misinterpretation when a group is
          recording the average of X and the maximum of Y, etc.
          within one data file makes such a capability more dangerous
          than useful.
    \item Note – It is easy to confuse, for example, a record of the
          maximum values of a set of variables, and the values of a set
          of variables associated with some specific point in time, such
          as when a specific variable reaches its maximum value.
          The values recorded to represent the maximum values of a
          set of variables are not necessarily temporally coincident
          (e.g. when recording atmospheric temperature and altitude,
          it is unlikely that the maximum temperature would be reached
          at the maximum altitude). A coincident data set from the
          point at which some variable reaches a maximum value can
          be acquired by setting a logging conditional to record only
          when that variable reaches its maximum value.
  \end{itemize}
\end{itemize}

\section {Mathematical Formulations}

\chapter {User's Guide}
\section {Implementation}
\section {Initialization}
\section {Routine Execution}
\section {Optional Execution}
\section {Configuration}

\chapter{Verification}

\end{document}
